\newpage
\chapter{Результаты измерений}

\section{Амплитудная модуляция}

\begin{table}[H]
    \small
    \renewcommand{\arraystretch}{1.3}
    \centering
    \begin{tabularx}{\textwidth}{|c|Y|Y|c|c|Y|}
        \hline
        $m$ & $U_{\text{нес}}$ ($12\text{ кГц}$), В & $U_{\text{бок}}$ ($11; 13\text{ кГц}$), В & $U_{\max}$, В & $U_{\min}$, В & Длина $L$, В \\ \hline
        0.3 & 1.0 & 0.15 & 1.30 & 0.70 & 0.60 \\ \hline
        0.4 & 1.0 & 0.20 & 1.40 & 0.60 & 0.80 \\ \hline
        0.5 & 1.0 & 0.25 & 1.50 & 0.50 & 1.00 \\ \hline
        0.6 & 1.0 & 0.30 & 1.60 & 0.40 & 1.20 \\ \hline
        0.7 & 1.0 & 0.35 & 1.70 & 0.30 & 1.40 \\ \hline
        0.8 & 1.0 & 0.40 & 1.80 & 0.20 & 1.60 \\ \hline
        0.9 & 1.0 & 0.45 & 1.90 & 0.10 & 1.80 \\ \hline
        1.0 & 1.0 & 0.50 & 2.00 & 0.00 & 2.00 \\ \hline
    \end{tabularx}
    \caption{Результаты исследования характеристик АМ-сигнала ($U_0 = 1.0\text{ В}$, $f_0 = 12.0\text{ кГц}$)}
    \label{tab:AM_measurements}
\end{table}

\section{Частотная модуляция}

\begin{table}[H]
    \small
    \renewcommand{\arraystretch}{1.3}
    \centering
    \begin{tabularx}{\textwidth}{|c|Y|Y|Y|Y|Y|Y|Y|}
        \hline
        $m$ & Несущая $f_0 = 12\text{ кГц}$, В & Боковые $12 \pm 1\text{ кГц}$, В & Боковые $12 \pm 2\text{ кГц}$, В & Боковые $12 \pm 3\text{ кГц}$, В & Боковые $12 \pm 4\text{ кГц}$, В & Изме\-ренная полоса $\Delta f$, кГц & Вид фазового портрета \\ \hline
        0.2 & 0.99 & 0.10 & 0.005 & — & — & 2.0 & Короткая дуга \\ \hline
        1.0 & 0.77 & 0.44 & 0.11 & 0.02 & — & 4.0 & Дуга \\ \hline
        1.8 & 0.34 & 0.58 & 0.31 & 0.10 & 0.02 & 6.0 & Дуга \\ \hline
        2.6 & 0.10 & 0.47 & 0.46 & 0.24 & 0.08 & 8.0 & Дуга \\ \hline
        3.4 & 0.36 & 0.18 & 0.47 & 0.37 & 0.19 & 10.0 & Замкнутое кольцо \\ \hline
        4.2 & 0.38 & 0.14 & 0.31 & 0.43 & 0.31 & 12.0 & Замкнутое кольцо \\ \hline
        5.0 & 0.18 & 0.33 & 0.05 & 0.36 & 0.39 & 14.0 & Замкнутое кольцо \\ \hline
    \end{tabularx}
    \caption{Результаты исследования характеристик ЧМ-сигнала ($U_0 = 1.0\text{ В}$, $f_0 = 12.0\text{ кГц}$)}
    \label{tab:ChM_measurements}
\end{table}

\section{Фазовая модуляция}

\begin{table}[H]
    \small
    \renewcommand{\arraystretch}{1.3}
    \centering
    \begin{tabularx}{\textwidth}{|l|Y|Y|}
        \hline
        Параметр & Режим 1 ($m = 2.0$) & Режим 2 ($m = 4.0$) \\ \hline
        Амплитуда несущей $U_0$ & 1.0 В & 1.0 В \\ \hline
        Характер огибающей & Постоянная ($U(t) = \text{const} = 1.0\text{ В}$) & Постоянная ($U(t) = \text{const} = 1.0\text{ В}$) \\ \hline
        Угловой сектор развертки $\Delta\theta = 2m$ & $4.0\text{ рад } (\approx 229.2^\circ)$ & $8.0\text{ рад } (\approx 458.4^\circ > 360^\circ)$ \\ \hline
        Вид фазового портрета & Разомкнутая дуга окружности & Полностью замкнутое кольцо \\ \hline
        Радиус траектории на IQ-плоскости & $R = 1.0\text{ В}$ & $R = 1.0\text{ В}$ \\ \hline
    \end{tabularx}
    \caption{Результаты исследования параметров ФМ-сигнала ($U_0 = 1.0\text{ В}$, $f_0 = 12.0\text{ кГц}$)}
    \label{tab:PM_measurements}
\end{table}

\section{Двухпозиционная фазовая манипуляция}

\begin{table}[H]
    \small
    \renewcommand{\arraystretch}{1.3}
    \centering
    \begin{tabularx}{\textwidth}{|l|Y|Y|}
        \hline
        Параметр & Значение / Характеристика & Примечание / Источник \\ \hline
        Позиционность манипуляции ($M$) & 2 & 2 состояния фазы ($0^\circ$ и $180^\circ$) \\ \hline
        Разрядность символа ($n = \log_2 M$) & 1 бит/символ & 1 символ кодирует ровно 1 бит \\ \hline
        Фазовый сдвиг при смене бита ($\Delta\varphi$) & $180^\circ$ & Инверсия фазы несущей \\ \hline
        Координаты точек созвездия $(I; Q)$ & $(-1.0\text{ В}; 0)$ и $(+1.0\text{ В}; 0)$ & Код «0» $\rightarrow -1\text{ В}$; Код «1» $\rightarrow +1\text{ В}$ \\ \hline
        Траектория перехода луча & Прямая линия по оси $I$ & Проходит через центр координат $(0;0)$ \\ \hline
        Символьная скорость ($R_s$) & 1.0 кБод & По длительности символа $T_s = 1\text{ мс}$ \\ \hline
        Информационная скорость ($R$) & 1.0 кбит/с & Зафиксировано на Табло Д ($R = R_s \cdot n$) \\ \hline
    \end{tabularx}
    \caption{Измеренные параметры сигнала ФМ-2}
    \label{tab:BPSK_measurements}
\end{table}

\section{Сравнение характеристик сигналов ФМ-2 и ФМ-4}

\begin{table}[H]
    \small
    \renewcommand{\arraystretch}{1.3}
    \centering
    \begin{tabularx}{\textwidth}{|l|Y|Y|Y|}
        \hline
        Параметр & ФМ-2 (BPSK) & ФМ-4 (QPSK) & Примечание \\ \hline
        Позиционность ($M$) & 2 & 4 & Число фазовых состояний \\ \hline
        Группировка бит ($n = \log_2 M$) & 1 бит & 2 бита & Коды 00, 01, 10, 11 \\ \hline
        Количество точек созвездия & 2 & 4 & По 1 в каждом квадранте \\ \hline
        Фазовые углы точек & $0^\circ, 180^\circ$ & $45^\circ$, $135^\circ$, $225^\circ$, $315^\circ$ & Симметрично осям $I$ и $Q$ \\ \hline
        Фазовый сдвиг & $180^\circ$ ($\pi$ рад) & $90^\circ$ ($\pi/2$ рад) & Шаг между соседними точками \\ \hline
        Символьная скорость ($R_s$) & 1.0 кБод & 1.0 кБод & Длительность символа $T_s = 1\text{ мс}$ \\ \hline
        Информационная скорость ($R$) & 1.0 кбит/с & 2.0 кбит/с & Увеличилась в 2 раза \\ \hline
        Ширина главного лепестка спектра & 2.0 кГц & 2.0 кГц & Занимаемая полоса не изменилась \\ \hline
    \end{tabularx}
    \caption{Сравнительные характеристики сигналов ФМ-2 и ФМ-4}
    \label{tab:PM_comparison}
\end{table}

\section{Двухпозиционная частотная манипуляция}

\begin{table}[H]
    \centering
    \small
    \renewcommand{\arraystretch}{1.3}
    \begin{tabularx}{\textwidth}{|l|Y|Y|}
        \hline
        Параметр & Значение & Примечание \\ \hline
        Центральная несущая ($f_0$) & 12.0 кГц & Базовая частота генератора \\ \hline
        Количество выраженных пиков в спектре & 2 пика & По одному на каждый логический уровень \\ \hline
        Частоты спектральных пиков ($f_1$ и $f_2$) & $9.0\text{ кГц}$ и $15.0\text{ кГц}$ & $f_1$ — для бита «0», $f_2$ — для бита «1» \\ \hline
        Частотный разнос между пиками ($\Delta F$) & $6.0\text{ кГц}$ & $\Delta F = |f_2 - f_1| = 2 \cdot \Delta f$ (девиация $\Delta f = 3.0\text{ кГц}$) \\ \hline
        Форма фазового портрета & Сплошное кольцо & Радиус $R = 1.0\text{ В}$ (непрерывное вращение) \\ \hline
        Информационная скорость ($R$) & 1.0 кбит/с & Табло Д \\ \hline
    \end{tabularx}
    \caption{Результаты измерения параметров сигнала ЧМ-2}
    \label{tab:BFSK_measurements}
\end{table}

\section{Четырехпозиционная частотная манипуляция}

\begin{table}[H]
    \centering
    \small
    \renewcommand{\arraystretch}{1.3}
    \begin{tabularx}{\textwidth}{|l|Y|Y|X|}
        \hline
        Параметр & ЧМ-2 & ЧМ-4 & Примечание \\ \hline
        Центральная несущая ($f_0$) & 12.0 кГц & 12.0 кГц & Опорная частота \\ \hline
        Девиация частоты ($\Delta f$) & 3.0 кГц & 4.5 кГц & За\-фик\-си\-ро\-вано на Табло Д \\ \hline
        Количество главных пиков в спектре & 2 пика & 4 пика & По одному на каждую комбинацию \\ \hline
        Частоты спектральных пиков & $9.0$ и $15.0\text{ кГц}$ & $7.5$; $10.5$; $13.5$; $16.5$ кГц & Сетка частот дибитов \\ \hline
        Частотный шаг между позициями ($\delta f$) & — & $3.0\text{ кГц}$ & Шаг между соседними пиками \\ \hline
        Полоса частот по крайним пикам & $6.0\text{ кГц}$ & $9.0\text{ кГц}$ & Полоса расширилась в 1.5 раза \\ \hline
        Информационная скорость ($R$) & 1.0 кбит/с & 2.0 кбит/с & Увеличилась в 2 раза (Табло Д) \\ \hline
        Спектральная эффективность ($R / \Delta F$) & $0.17\text{ бит}/(\text{с}\cdot\text{Гц})$ & $0.22\text{ бит}/(\text{с}\cdot\text{Гц})$ & Нез\-на\-чи\-тель\-ный прирост \\ \hline
    \end{tabularx}
    \caption{Сравнительные характеристики сигналов ЧМ-2 и ЧМ-4}
    \label{tab:QFSK_measurements}
\end{table}

\section{Амплитудно-фазовая манипуляция (АФМ-16)}

\begin{table}[H]
    \centering
    \small
    \renewcommand{\arraystretch}{1.3}
    \begin{tabularx}{\textwidth}{|l|Y|Y|}
        \hline
        Параметр & Значение / Характеристика & Примечание / Источник \\ \hline
        Позиционность манипуляции ($M$) & 16 & Число состояний радиосигнала \\ \hline
        Разрядность символа ($n = \log_2 M$) & 4 бита & 1 символ переносит сразу 4 бита \\ \hline
        Геометрия созвездия на IQ & 2 концентрических кольца & Структура стандарта спутникового ТВ DVB-S2 \\ \hline
        Внутреннее кольцо ($R_1$) & 4 точки (шаг $90^\circ$) & Расположены под углами $45^\circ$, $135^\circ$, $225^\circ$, $315^\circ$ \\ \hline
        Внешнее кольцо ($R_2$) & 12 точек (шаг $30^\circ$) & Равномерно с шагом $\pi/6$ ($0^\circ$, $30^\circ$, $60^\circ \dots$) \\ \hline
        Соотношение радиусов ($R_2 / R_1$) & $\approx 2.5 \dots 2.8$ & Обеспечивает максимальную помехоустойчивость \\ \hline
        Характер огибающей (Панель А) & Переменная (2 уровня) & Изменяется скачкообразно между $R_1$ и $R_2$ \\ \hline
        Символьная скорость ($R_s$) & 1.0 кБод & Длительность символа $T_s = 1.0\text{ мс}$ \\ \hline
        Информационная скорость ($R$) & 4.0 кбит/с & В 4 раза выше ФМ-2 (Табло Д) \\ \hline
    \end{tabularx}
    \caption{Результаты измерения параметров сигнала АФМ-16}
    \label{tab:APSK-16_measurements}
\end{table}

\section{Амплитудно-фазовая манипуляция (АФМ-32)}

\begin{table}[H]
    \centering
    \small
    \renewcommand{\arraystretch}{1.3}
    \begin{tabularx}{\textwidth}{|l|Y|Y|}
        \hline
        Параметр & Значение / Характеристика & Примечание / Источник \\ \hline
        Позиционность манипуляции ($M$) & 32 & Число состояний радиосигнала \\ \hline
        Разрядность символа ($n = \log_2 M$) & 5 бит & 1 символ переносит сразу 5 бит \\ \hline
        Геометрия созвездия на IQ & 3 концентрических кольца & Высокая плотность \\ \hline
        Внутреннее кольцо ($R_1$) & 4 точки (шаг $90^\circ$) & Углы $45^\circ$, $135^\circ$, $225^\circ$, $315^\circ$ \\ \hline
        Среднее кольцо ($R_2$) & 12 точек (шаг $30^\circ$) & Равномерно с угловым шагом $\pi/6$ \\ \hline
        Внешнее кольцо ($R_3$) & 16 точек (шаг $22.5^\circ$) & Равномерно с угловым шагом $\pi/8$ \\ \hline
        Количество уровней амплитуды (Панель А) & 3 уровня & Соответствуют радиусам $R_1, R_2, R_3$ \\ \hline
        Символьная скорость ($R_s$) & 1.0 кБод & Постоянна ($T_s = 1.0\text{ мс}$) \\ \hline
        Скорость на Табло Д & 5.0 кбит/с & Информационная скорость ($R$) \\ \hline
    \end{tabularx}
    \caption{Результаты измерения параметров сигнала АФМ-32}
    \label{tab:APSK-32_measurements}
\end{table}

\section{Квадратурная амплитудная манипуляция (КАМ-16)}

\begin{table}[H]
    \centering
    \small
    \renewcommand{\arraystretch}{1.5}
    \begin{tabularx}{\textwidth}{|l|Y|Y|}
        \hline
        Параметр & Значение / Характеристика & Примечание / Источник \\ \hline
        Позиционность модуляции ($M$) & 16 & $M = 2^4$ состояний \\ \hline
        Разрядность символа ($n = \log_2 M$) & 4 бита & 1 символ кодирует 4 бита \\ \hline
        Геометрия созвездия на IQ (Панель Г) & Квадратная матрица $4 \times 4$ & 16 неподвижных точек на сетке \\ \hline
        Дискретные уровни осей $I$ и $Q$ & $\pm 0.33\text{ В}; \pm 1.0\text{ В}$ & По формуле: $[-3, -1, +1, +3] \cdot \frac{U_0}{3}$ \\ \hline
        Максимальная пиковая амплитуда $U_m$ & $1.414\text{ В}$ & $U_m = \sqrt{1.0^2 + 1.0^2} = \sqrt{2}\text{ В}$ \\ \hline
        Информационная скорость ($R$) & 4.0 кбит/с & В 4 раза выше ФМ-2 \\ \hline
        Ширина главного лепестка спектра & $\approx 2.0\text{ кГц}$ & Как у ФМ-2 / АФМ-16 ($R_s = 1.0\text{ кБод}$) \\ \hline
    \end{tabularx}
    \caption{Результаты исследования характеристик сигнала КАМ-16}
    \label{tab:QAM-16_measurements}
\end{table}

\section{Квадратурная амплитудная манипуляция (КАМ-32)}

\begin{table}[H]
    \centering
    \small
    \renewcommand{\arraystretch}{1.5}
    \begin{tabularx}{\textwidth}{|l|Y|Y|}
        \hline
        Параметр & Значение / Характеристика & Примечание / Источник \\ \hline
        Позиционность модуляции ($M$) & 32 & $M = 2^5$ состояний \\ \hline
        Разрядность символа ($n = \log_2 M$) & 5 бит & 1 символ кодирует сразу 5 бит \\ \hline
        Геометрия созвездия на IQ & Созвездие-крест, 32 точки & Сетка $6 \times 6$ с 4 вырезанными углами \\ \hline
        Базовая координатная сетка & $I, Q \in [-5, -3, -1, +1, +3, +5] \cdot \frac{U_0}{4.5}$ & 6 дискретных уровней по каждой оси \\ \hline
        Координаты самых удаленных точек & $(\pm 5, \pm 3)$ и $(\pm 3, \pm 5)$ & В долях базового шага сетки \\ \hline
        Максимальная амплитуда $U_m$ & $1.296\text{ В} \approx 1.30\text{ В}$ & $U_m = \frac{\sqrt{5^2 + 3^2}}{4.5} = \frac{\sqrt{34}}{4.5}$ \\ \hline
        Информационная скорость ($R$) & 5.0 кбит/с & В 5 раз выше ФМ-2 \\ \hline
        Ширина главного лепестка спектра & $\approx 2.0\text{ кГц}$ & Как у АФМ-32 ($R_s = 1.0\text{ кБод}$) \\ \hline
    \end{tabularx}
    \caption{Результаты исследования характеристик сигнала КАМ-32}
    \label{tab:QAM-32_measurements}
\end{table}

